% !TEX program = xelatex
\documentclass[10pt,letterpaper,twocolumn]{article}

\usepackage[top=0.85in, bottom=0.85in, left=0.75in, right=0.75in]{geometry}
\usepackage{fontspec}
\usepackage{unicode-math}
\setmainfont{Latin Modern Roman}[
  Extension=.otf,
  UprightFont=lmroman10-regular,
  BoldFont=lmroman10-bold,
  ItalicFont=lmroman10-italic,
  BoldItalicFont=lmroman10-bolditalic,
]
\setsansfont{Latin Modern Sans}[
  Extension=.otf,
  UprightFont=lmsans10-regular,
  BoldFont=lmsans10-bold,
  ItalicFont=lmsans10-oblique,
  BoldItalicFont=lmsans10-boldoblique,
]
\newfontfamily\acbold{ywft-processing-bold}[Path=../../system/public/type/webfonts/,Extension=.ttf]
\newfontfamily\aclight{ywft-processing-light}[Path=../../system/public/type/webfonts/,Extension=.ttf]
\setmonofont{Latin Modern Mono}[
  Extension=.otf,
  UprightFont=lmmono10-regular,
  ItalicFont=lmmono10-italic,
  Scale=0.85,
]

\usepackage{xcolor}
\usepackage{titlesec}
\usepackage{enumitem}
\usepackage{booktabs}
\usepackage{fancyhdr}
\usepackage{graphicx}
\graphicspath{{figures/}{../../papers/arxiv-ac/figures/}}
\usepackage{ragged2e}
\usepackage{microtype}
\usepackage{eso-pic}
\usepackage{tikz}
\usepackage{hyperref}
\usepackage{xurl}

\definecolor{acpink}{RGB}{180,72,135}
\definecolor{acpurple}{RGB}{120,80,180}
\definecolor{acdark}{RGB}{64,56,74}
\definecolor{acgray}{RGB}{119,119,119}
\definecolor{palspink}{RGB}{205,92,155} % #cd5c9b — canonical purple-pals.svg fill

% === PALS SIDE-STAMPS (mirrors bios.mjs side-stamp pattern + ac-paper-layout.sty) ===
\AddToShipoutPictureBG{%
  \begin{tikzpicture}[remember picture, overlay]
    \pgfmathsetmacro{\palswiggle}{7*sin(\thepage * 67)}%
    \node[opacity=0.10, anchor=center, rotate={90 + \palswiggle}]
      at ([xshift=-1.0cm, yshift=-0.33\paperheight]current page.north west)
      {
\includegraphics[width=13cm]{palsmark}};
    \node[opacity=0.10, anchor=center, rotate={-90 + \palswiggle}]
      at ([xshift=1.0cm, yshift=-0.33\paperheight]current page.north east)
      {
\includegraphics[width=13cm]{palsmark}};
  \end{tikzpicture}%
}

\hypersetup{colorlinks=true,linkcolor=acpurple,urlcolor=acpurple,citecolor=acpurple,
  pdftitle={The Pals Mark: En historie om Aesthetic.Computer-logoet}}

\titleformat{\section}{\normalfont\bfseries\normalsize\uppercase}{\thesection.}{0.5em}{}
\titlespacing{\section}{0pt}{1.2em}{0.3em}
\titleformat{\subsection}{\normalfont\bfseries\small}{\thesubsection}{0.5em}{}
\titlespacing{\subsection}{0pt}{0.8em}{0.2em}

\pagestyle{fancy}\fancyhf{}
\renewcommand{\headrulewidth}{0pt}
\fancyhead[C]{\footnotesize\color{acpink}\textit{The Pals Mark --- \ac{}}}
\fancyfoot[C]{\footnotesize\thepage}

\setlist[itemize]{nosep, leftmargin=1.2em, itemsep=0.15em}
\setlist[enumerate]{nosep, leftmargin=1.4em}
\setlength{\columnsep}{1.8em}
\setlength{\parindent}{1em}
\setlength{\parskip}{0.3em}

\newcommand{\ac}{\textsc{Aesthetic}{\color{acpink}.}\textsc{Computer}}
\newcommand{\acdot}{{\color{acpink}.}}
\newcommand{\pals}{\textsc{Pals}}
% commit link: \commit{shorthash}{message} -> links the hash to GitHub
\newcommand{\commit}[2]{\href{https://github.com/whistlegraph/aesthetic-computer/commit/#1}{\texttt{#1}}, \emph{``#2''}}
% repo file link
\newcommand{\repofile}[1]{\href{https://github.com/whistlegraph/aesthetic-computer/blob/main/#1}{\texttt{#1}}}

% Cover language bar. The same paper is published in six languages; each
% edition links to the others by canonical permalink. \klanghere marks the
% current language; \klang links a sibling edition by its URL suffix.
\newcommand{\klang}[2]{\href{https://papers.aesthetic.computer/pals-mark-26-arxiv#2.pdf}{\color{acgray}#1}}
\newcommand{\klanghere}[1]{{\bfseries\color{acpink}#1}}
\newcommand{\klangsep}{{\color{acgray!60}\,\textperiodcentered\,}}

\begin{document}

\twocolumn[{%
\begin{center}
{\acbold\fontsize{27pt}{31pt}\selectfont\color{acdark} The \pals{} Mark}\par
\vspace{0.5em}
\includegraphics[width=0.6\textwidth]{figures/cover}\par
\vspace{0.4em}
{\itshape\fontsize{14.5pt}{17pt}\selectfont\color{acpink} En historie om Aesthetic{\acdot}Computer-logoet}\par
\vspace{0.5em}
{\small\color{acgray} \href{https://prompt.ac/@jeffrey}{@jeffrey} \,$\cdot$\, Aesthetic{\acdot}Computer \,$\cdot$\, logo af Niki Stebbins \,$\cdot$\, ORCID \href{https://orcid.org/0009-0007-4460-4913}{0009-0007-4460-4913}}\par
\vspace{0.45em}
{\footnotesize\sffamily\klang{EN}{}\klangsep\klanghere{DA}\klangsep\klang{ES}{-es}\klangsep\klang{ZH}{-zh}\klangsep\klang{JA}{-ja}\klangsep\klang{RU}{-ru}}\par
\vspace{0.55em}
{\color{acpink}\rule{\textwidth}{1.2pt}}
\vspace{0.5em}
\end{center}

\begin{quote}
\small\noindent\textbf{Resumé.}
\pals{} er logoet for \ac{}: to håndtegnede figurer, der holder i hånd med
armene løftet. Denne artikel fortæller, hvordan mærket blev til --- bestilt
hos kunstneren Niki Stebbins og tegnet, hen over sommeren 2023,
\emph{inde} i \ac{} med platformens eget \texttt{wgr}-whistlegraph-værktøj ---
og hvordan det derefter blev infrastruktur: en roterende
\texttt{pals.aesthetic.computer}-``logo-griber'', der serverer mærkets mange
positurer på tværs af avatarer, artikler, native OS og nettet, hvor hver
enkelt er et link tilbage til projektet.
\end{quote}
\vspace{0.5em}
}]

\section{Hvad \ac{} er}

\ac{} er en mobil-først runtime og et socialt netværk for kreativ
databehandling --- en enkelt URL-linje, der samtidig fungerer som en
musikinstrument-lignende grænseflade, hvor huskbare kommandoer fremkalder
interaktive programmer kaldet ``pieces''. Det understøtter både JavaScript
(\texttt{.mjs}) og en hjemmedyrket Lisp-dialekt, \textbf{KidLisp} (118
indbyggede funktioner). Hele systemet taler til pieces gennem ét kerne-API
(\texttt{disk.mjs}, $\sim$600\,KB).

\textbf{Hvornår det gik online.} Domænet \texttt{aesthetic.computer} har
været holdt siden \textbf{2013-12-12}. Platformen, som den findes i dag,
blev undfanget ud af \emph{No Paint} (2020), udviklingen begyndte i
\textbf{2021}, og den har været deployet kontinuerligt fra det nuværende
repository siden \textbf{december 2022} ($\sim$20{,}300 commits).

\textbf{Platformsstatistik} (live, \today):
\begin{itemize}
  \item $\sim$400 indbyggede pieces i repositoriet (383 \texttt{.mjs} + 20
        \texttt{.lisp}); 265 brugerudgivne værker.
  \item \textbf{17{,}538} KidLisp-programmer lagret i produktion --- et
        voksende korpus af fællesskabsforfattet generativ kunst.
  \item ATProto/Bluesky-federering; per-handle-samlinger af records
        (malerier, stemninger, pieces, KidLisp) på \texttt{at://}-netværket.
\end{itemize}

\textbf{Hvem det tjener, og hvem der kender mærket.} \ac{}'s brugere er et
fællesskab af kreative kodere, kunstnere og ikke-tekniske ``Pals'', der
udgiver øjeblikkeligt under \texttt{@handle}-pieces. Platformen engagerer et
bredt læserskare på tværs af kunst-og-teknologi-feltet, og \ac{}-værker er
kommet ind i institutionelle samlinger. Selve
\pals{}-mærket rider på hver eneste af disse overflader: boot-skærmen,
artiklerne, de sociale avatarer, native OS. Dets begrebsmæssige slægtslinje
i det radikale computermusik-fællesskab \emph{Goodiepal \& Pals} (se
slægtsafsnittet nedenfor) giver navnet en anerkendt historie i medie\-kunsten
ud over \ac{}'s eget publikum.

\section{Hvad mærket forestiller}

\pals{}-logoet er \textbf{to optegnede menneskefigurer, der holder i hånd
med armene løftet} --- som to venner, der går eller danser sammen ---
gengivet i en løs, organisk, håndtegnet streg. Den kanoniske statiske
vektor (\texttt{system/public/purple-pals.svg}) er pink/magenta,
\texttt{\#cd5c9b}. Det er \emph{ikke} ``et par øjne'', en tidligere
fejllæsning, som siden er blevet rettet i artikelstakken.

Afgørende er, at mærket blev lavet \emph{med} platformen, ikke blot
\emph{til} den: Niki tegnede logo-udforskningerne med \texttt{wgr}, det
whistlegraph-renderende tegne-piece, der kører inde i selve \ac{}. Logoet
er derfor både en identitet og en demonstration af værktøjet.

\section{Hvordan mærket blev til}

\subsection{Først et stempel}
Det begynder i sommeren \textbf{2020} med en lille gave. Niki Stebbins
sender en tegning på mail --- en lille lilla stjerne --- og Jeffrey lægger
den ind i spillet som et stempel under handlet \texttt{@niki} og beder om
mere. De næste to år er samarbejdet en løs, vedvarende \emph{tegnesession}:
stempler, et tegnespil, whistlegraph-kruseduller. I \textbf{2022} er det
betalt designarbejde. Intet af det er logoet endnu, men det er den jord,
logoet vokser af --- en vane, hvor to mennesker tegner sammen inde i
værktøjet.

\subsection{Designet genåbnes (juni 2023)}
Selve logoet begynder med utilfredshed. I \textbf{juni 2023} skriver Niki,
at hun ``ikke er 100\% tilfreds med det logo, jeg efterlod dig'', og genåbner
det. Hun tænker højt om farve --- ``som lilla eller grøn\dots{} ikke blå, der
er for mange blå app-logoer'' --- og om form: skulle mærket være ``som en
ramme, kun en lukket figur'', noget man kunne arbejde \emph{inde} i? Hun
sender en skitsebogsside med kandidater og udpeger en favorit --- ``jeg
kan bedst lide den med cirklen om'' --- og bekymrer sig så straks for, at
det er ``ret lig are.nas logo'', og fortsætter. Den nøgleinstruktion, hun
giver sig selv, er en proportionskorrektion: skub det ``så det ligner mindre
stjerner og mere figurer''. En side af studierne får sin egen dom ---
``ligner søstjerner lol'' --- og for at rette det forlader hun skærmen helt
og binder og fotograferer længder af \textbf{reb} i en park for at forstå,
hvordan to sammenkoblede kroppe faktisk falder.

\subsection{De to figurer bygges (juli--august 2023)}
Hen over sommeren opløser mærket sig til \emph{to} figurer frem for én.
Niki køber et brugt \textbf{Kodak EasyShare}-digitalkamera specifikt for at
fotografere logoet, dimensionerer en version til ``den nøjagtige størrelse
af VS Marketplace-logoet'' og spørger, om Jeffrey vil have ``flere
logotegninger i a.c til pfp'er på sociale medier og app-store-tingen''.
Repositoriet stemmer overens med beskederne: på
\commit{5bcefc5a9}{better defaults in \texttt{wgr} for niki's logo efforts}
bliver selve tegneværktøjet tunet til hende, og ved
\commit{a0009f5d0}{add new vscode extension logo from niki} lander et
afsendt aktiv. Sidst i august melder hun ``varianterne af de \textbf{to
logoer}'' færdige og beskriver gesten præcist --- ``den højre figur trak i
den venstre\dots{} den venstre figurs arm er højere''. Mærket er nu det,
det stadig er: to mennesker, arme oppe, den ene trækker den anden med sig.
Den \textbf{2023-08-31} vælges et Instagram-profilbillede.

\subsection{Softice, og live-gang (september 2023)}
Afslutningen er den mærkeligste og bedste del. For at fotografere mærket ude
i verden tegner Niki det i fuld størrelse på fortovet og forsøger at skyde
det med is --- men isen fra boldbanen kommer i bægre, så hun tager en omvej
til en McDonald's efter en vaffelis og får billedet. Den \textbf{2023-09-09}
``sendte hun logoerne til a.c-mailen'' under emnet \emph{!c3 cr34m l0go\$}
(brødtekst: ``ice cream logos 8)''). Samme uge holder mærket op med at være
et billede og bliver infrastruktur: \commit{0aba1a378}{add dynamic logo
generator} lander, og Jeffrey melder tilbage det etlinjede resultat ---
``\^{}\ vil nu returnere et tilfældigt logo hver gang''. Logoet blev tegnet
\emph{med} platformen (i \texttt{wgr}-whistlegraph-værktøjet) og derefter
serveret \emph{af} den. Det var aldrig kun et billede; det var altid også et
program.

\section{Endpointet: \texttt{pals.aesthetic.computer}}

Mærket blev infrastruktur i september 2023, da \ac{} rejste en dedikeret
logo-tjeneste og gav den et huskbart hjem.

\begin{itemize}
  \item \textbf{2023-09-05} --- \commit{0aba1a378}{add dynamic logo
        generator}. Backend-modulet \repofile{system/backend/logo.mjs}
        er stemplet \texttt{23.09.05.22.31}.
  \item \textbf{2023-09-12} --- \commit{f1a4f272c}{logo -> pals.aesthetic}
        --- ruten omdøbes til \pals{}-identiteten.
  \item \textbf{2023-09-13} --- \commit{12e348035}{pals.aesthetic.computer
        url fix} --- endpointet er live og adresserbart på sit eget
        subdomæne.
\end{itemize}

\noindent Så \href{https://pals.aesthetic.computer}{\texttt{pals.aesthetic.computer}}
\textbf{gik online 2023-09-12--13}, tre dage efter at den dynamiske generator
landede.

\subsection{Hvordan det virker som en ``logo-griber''}

Endpointet er en lillebitte roterende proxy, ikke et enkelt fast billede.
Hver forespørgsel returnerer en anden \pals{}-positur:

\begin{enumerate}
  \item \href{https://pals.aesthetic.computer}{\texttt{pals.aesthetic.computer}}
        omdirigerer (via \repofile{system/netlify.toml}) til
        \href{https://aesthetic.computer/api/logo}{\texttt{aesthetic.computer/api/logo}}.
  \item \repofile{system/backend/logo.mjs} holder en liste over
        \pals{}-\emph{malerier} på CDN-bucketen
        \texttt{pals-aesthetic-computer} og returnerer et tilfældigt et
        ved hvert kald (og undgår en umiddelbar gentagelse).
  \item Funktionen indholds-forhandler: en browser får en fuldflade
        HTML-side på en tilfældig lilla, som man kan trykke på for at cykle
        til den næste positur (den sender \texttt{?previousLogo=} med, så
        den aldrig gentager); en \texttt{curl}/\texttt{<img>}/server-forespørgsel
        får den rå \texttt{.png} direkte. Den samme URL er dermed både en
        fremviser og et hot-linkbart aktiv.
\end{enumerate}

\subsection{De mange positurer af \pals{}}

Griberens positursæt er et galleri af \texttt{wgr}-malerier fra
midten af 2023 (serien \texttt{painting-2023.7.29}--\texttt{2023.8.21}).
Fordi logoet er en \emph{samling af tegnede positurer} snarere end én
frossen glyf, læses varemærket som levende: hver sideindlæsning,
avatar-opdatering eller favicon-hentning kan vise de to figurer i en let
forskellig stilling. Det fulde sæt af tolv live positurer er gengivet i
Appendiks~A.

\section{Brug på tværs af egenskaber}

Fordi enhver egenskab kan hot-linke
\href{https://pals.aesthetic.computer}{\texttt{pals.aesthetic.computer}}
(eller indlejre den statiske
\repofile{system/public/purple-pals.svg}), tråder ét enkelt mærke sig
gennem hele økosystemet, og hver forekomst er et link tilbage til projektet:

\begin{itemize}
  \item \textbf{Sociale avatarer \& Gravatar} --- det roterende endpoint
        leverer et profilbillede, der stille skifter positur, så den samme
        identitet fremtræder på tværs af tjenester, mens den altid opløses
        til \ac{}.
  \item \textbf{Web \& marketing} --- favicons, doc-headers og indlejrede
        \texttt{<img>}-tags trækker de live \pals{}; et tryk cykler
        positurer og forvandler selve logoet til en lille interaktiv
        teaser for platformen.
  \item \textbf{Artikler} --- \pals{} rider i marginerne af ethvert
        \ac{}-arbejdsudkast som halvgennemsigtige, side-vrikkende
        side-stempler (se denne side) og optræder i tryk\-udgavernes
        kolofoner.
  \item \textbf{Native OS \& desktop} --- \pals{} tjener som boot-branding
        (Plymouth-splash) i AC Native OS og som vektor-tray-ikonet i AC
        Electron og MenuBand --- accent-tonet og badge-bar.
  \item \textbf{Tryk (afventer)} --- i maj 2025 kontaktede Jeffrey Casey
        Rubber Stamps for at skære fysiske \pals{}-stempler i tre
        størrelser; ordren gik i stå og blev aldrig fuldført.
\end{itemize}

\section{Hvorfor ``\pals{}''? Goodiepal-slægtslinjen}

Navnet er ikke udsmykning; det er en strukturel påstand, og dets vægt
kommer udefra \ac{}. Som argumenteret i \ac{}-artiklen
\emph{Radical Computer Art: Goodiepalian Approaches in Aesthetic
Computer} (\repofile{papers/arxiv-goodiepal/goodiepal.tex}) og den
platformsartikel, den citerer, nedstammer ``\pals{}'' fra \textbf{Goodiepal
\& Pals} --- den løse affinitetsgruppe omkring den dansk-færøske musiker og
Radical Computer Music-teoretiker Goodiepal (Gaeoudji SYGNOK, f.~1974).

\pals{}-modellen er en organisationspolitik: en gruppe uden medlemskab,
uden institutionel opbakning og uden formel struktur, holdt sammen af delt
praksis og gensidig anerkendelse snarere end tilhørsforhold. Goodiepal
underviste i elektronisk musiks historie og æstetik ved Det Kongelige
Danske Musikkonservatorium og forlod det i 2008 efter en konflikt beskrevet
som at have ``den forkerte slags venner'', efter at have erklæret
``intellektuel krig mod dumheden i moderne computermusik og mediekunst''.
Jeffrey turnerede med Goodiepal \& Pals rundt i Europa i 2018, og de
værktøjer, der blev bygget til de optrædener, er direkte forfædre til
\ac{}'s arkitektur.

Så logoet bærer den mediekunstneriske politiske historie ind i varemærket.
Som Goodiepal-artiklen udtrykker det, er det at navngive logoet
``\pals{}'' --- og at give partituret titlen \emph{Score for
Aesthetic{\color{acpink}.}Computer \& Pals} --- ``ikke branding; det er en
strukturel påstand: platformen tilhører sit fællesskab af praktikere, ikke
dens udvikler eller dens investorer''. To figurer, der holder i hånd, lavet
\emph{inde} i værktøjet, er argumentet gengivet som et mærke: \ac{} er
\emph{for} enhver, der dukker op --- for \pals{}.

\section{Ophav}

Jeffrey hyrede \textbf{Niki Stebbins} til at designe \ac{}-logoet under sin
ledelse. Hen over 2022--2023 udviklede hun mærket inde i platformens eget
\texttt{wgr}-tegneværktøj, arbejdede sig gennem skitsebogs-, reb- og
softice-udforskningerne dokumenteret ovenfor og leverede to-figurers-mærket,
der afsendes som \texttt{purple-pals.svg}. Den roterende
\texttt{pals.aesthetic.computer}-griber gik live 2023-09-12--13; den
kanoniske statiske vektor optræder første gang 2024-01-31
(\commit{bf29eadf8}{host purple-pals graphic for docs}); og fra 2026-02
centrerer repositoriet og artiklerne eksplicit ``\pals{}'' som
platformsidentitet (\commit{48dd8923e}{rename score to ``Score for
Aesthetic.Computer \& Pals''}). Mærket er et bestilt værk: forfattet af
Niki, ledet af Jeffrey og ejet af \ac{}.

\vspace{0.6em}
{\footnotesize\color{acgray}\noindent Kilder: repositoriets commit-historik;
\repofile{system/backend/logo.mjs};
\repofile{system/netlify/functions/logo.js};
\repofile{system/netlify.toml};
\repofile{papers/arxiv-goodiepal/goodiepal.tex};
\repofile{system/public/purple-pals.svg}.\par}

% ============================ APPENDICES ============================
\onecolumn
\appendix

\noindent\emph{Bemærkning om datoer.} Hvert visuelt aktiv nedenfor er
tekstet med den dato, det blev lavet: \texttt{wgr}-positurer med det
tidsstempel, der er indkodet i maleriets slug; monorepo-aktiver med deres
første-commit-dato (git \texttt{--diff-filter=A}); live-optagelser med den
dato, de blev fotograferet fra produktion.

\section*{Appendiks A --- De mange positurer (live \texttt{pals.aesthetic.computer}-sæt)}
\addcontentsline{toc}{section}{Appendiks A --- Positurer}
\noindent De tolv \texttt{wgr}-malerier, som logo-griberen roterer igennem,
trukket live fra CDN'en \texttt{pals-aesthetic-computer}. Samme
to-figurers-motiv, hver med en anden stilling og farvebehandling.
\vspace{1.2em}

\begin{minipage}[t]{0.242\textwidth}\centering
  
\includegraphics[width=\linewidth,height=\linewidth,keepaspectratio]{figures/poses/small/pose-1.png}\\[2pt]
  {\scriptsize\textbf{pose 1}}\\[-1pt]
  {\scriptsize\color{acgray}2023-08-01}
\end{minipage}\hfill\begin{minipage}[t]{0.242\textwidth}\centering
  \includegraphics[width=\linewidth,height=\linewidth,keepaspectratio]{figures/poses/small/pose-2.png}\\[2pt]
  {\scriptsize\textbf{pose 2}}\\[-1pt]
  {\scriptsize\color{acgray}2023-08-21}
\end{minipage}\hfill\begin{minipage}[t]{0.242\textwidth}\centering
  
\includegraphics[width=\linewidth,height=\linewidth,keepaspectratio]{figures/poses/small/pose-3.png}\\[2pt]
  {\scriptsize\textbf{pose 3}}\\[-1pt]
  {\scriptsize\color{acgray}2023-08-21}
\end{minipage}\hfill\begin{minipage}[t]{0.242\textwidth}\centering
  \includegraphics[width=\linewidth,height=\linewidth,keepaspectratio]{figures/poses/small/pose-4.png}\\[2pt]
  {\scriptsize\textbf{pose 4}}\\[-1pt]
  {\scriptsize\color{acgray}2023-08-01}
\end{minipage}

\vspace{0.8em}

\begin{minipage}[t]{0.242\textwidth}\centering
  \includegraphics[width=\linewidth,height=\linewidth,keepaspectratio]{figures/poses/small/pose-5.png}\\[2pt]
  {\scriptsize\textbf{pose 5}}\\[-1pt]
  {\scriptsize\color{acgray}2023-08-21}
\end{minipage}\hfill\begin{minipage}[t]{0.242\textwidth}\centering
  \includegraphics[width=\linewidth,height=\linewidth,keepaspectratio]{figures/poses/small/pose-6.png}\\[2pt]
  {\scriptsize\textbf{pose 6}}\\[-1pt]
  {\scriptsize\color{acgray}2023-08-01}
\end{minipage}\hfill\begin{minipage}[t]{0.242\textwidth}\centering
  \includegraphics[width=\linewidth,height=\linewidth,keepaspectratio]{figures/poses/small/pose-7.png}\\[2pt]
  {\scriptsize\textbf{pose 7}}\\[-1pt]
  {\scriptsize\color{acgray}2023-08-21}
\end{minipage}\hfill\begin{minipage}[t]{0.242\textwidth}\centering
  \includegraphics[width=\linewidth,height=\linewidth,keepaspectratio]{figures/poses/small/pose-8.png}\\[2pt]
  {\scriptsize\textbf{pose 8}}\\[-1pt]
  {\scriptsize\color{acgray}2023-08-21}
\end{minipage}

\vspace{0.8em}

\begin{minipage}[t]{0.242\textwidth}\centering
  \includegraphics[width=\linewidth,height=\linewidth,keepaspectratio]{figures/poses/small/pose-9.png}\\[2pt]
  {\scriptsize\textbf{pose 9}}\\[-1pt]
  {\scriptsize\color{acgray}2023-08-01}
\end{minipage}\hfill\begin{minipage}[t]{0.242\textwidth}\centering
  \includegraphics[width=\linewidth,height=\linewidth,keepaspectratio]{figures/poses/small/pose-10.png}\\[2pt]
  {\scriptsize\textbf{pose 10}}\\[-1pt]
  {\scriptsize\color{acgray}2023-08-01}
\end{minipage}\hfill\begin{minipage}[t]{0.242\textwidth}\centering
  \includegraphics[width=\linewidth,height=\linewidth,keepaspectratio]{figures/poses/small/pose-11.png}\\[2pt]
  {\scriptsize\textbf{pose 11}}\\[-1pt]
  {\scriptsize\color{acgray}2023-07-31}
\end{minipage}\hfill\begin{minipage}[t]{0.242\textwidth}\centering
  
\includegraphics[width=\linewidth,height=\linewidth,keepaspectratio]{figures/poses/small/pose-12.png}\\[2pt]
  {\scriptsize\textbf{pose 12}}\\[-1pt]
  {\scriptsize\color{acgray}2023-07-29}
\end{minipage}

\vspace{0.8em}




\newpage
\section*{Appendiks B --- Mærket på tværs af monorepoet}
\addcontentsline{toc}{section}{Appendiks B --- Monorepo-aktiver}
\noindent De samme \pals{} tegnet én gang og genskindet overalt: den
kanoniske vektor og dens efterkommere på tværs af kodebasen ---
editor-udvidelse, native-OS boot-splash, desktop-tray, pop-covers,
faktura-stempel og keymaps-farvesættet. Hver flise er tekstet med sin
første-commit-dato.
\vspace{1.2em}

\begin{minipage}[t]{0.194\textwidth}\centering
  
\includegraphics[width=\linewidth,height=\linewidth,keepaspectratio]{figures/assets/c01-purple-pals.png}\\[2pt]
  {\scriptsize\textbf{purple-pals.svg}}\\[-1pt]
  {\scriptsize\color{acgray}2024-01-31}
\end{minipage}\hfill\begin{minipage}[t]{0.194\textwidth}\centering
  
\includegraphics[width=\linewidth,height=\linewidth,keepaspectratio]{figures/assets/c02-vscode.png}\\[2pt]
  {\scriptsize\textbf{vscode-extension}}\\[-1pt]
  {\scriptsize\color{acgray}2024-01-31}
\end{minipage}\hfill\begin{minipage}[t]{0.194\textwidth}\centering
  
\includegraphics[width=\linewidth,height=\linewidth,keepaspectratio]{figures/assets/c03-tvos.png}\\[2pt]
  {\scriptsize\textbf{apple/tvos-tapes}}\\[-1pt]
  {\scriptsize\color{acgray}2024-01-31}
\end{minipage}\hfill\begin{minipage}[t]{0.194\textwidth}\centering
  
\includegraphics[width=\linewidth,height=\linewidth,keepaspectratio]{figures/assets/c04-plymouth.png}\\[2pt]
  {\scriptsize\textbf{AC-OS boot splash}}\\[-1pt]
  {\scriptsize\color{acgray}2026-02-21}
\end{minipage}\hfill\begin{minipage}[t]{0.194\textwidth}\centering
  
\includegraphics[width=\linewidth,height=\linewidth,keepaspectratio]{figures/assets/c05-tray.png}\\[2pt]
  {\scriptsize\textbf{Electron tray}}\\[-1pt]
  {\scriptsize\color{acgray}2026-05-20}
\end{minipage}

\vspace{0.8em}

\begin{minipage}[t]{0.194\textwidth}\centering
  
\includegraphics[width=\linewidth,height=\linewidth,keepaspectratio]{figures/assets/c14-stamp.png}\\[2pt]
  {\scriptsize\textbf{invoice stamp}}\\[-1pt]
  {\scriptsize\color{acgray}2026-03-06}
\end{minipage}\hfill\begin{minipage}[t]{0.194\textwidth}\centering
  
\includegraphics[width=\linewidth,height=\linewidth,keepaspectratio]{figures/assets/c12-gray.png}\\[2pt]
  {\scriptsize\textbf{grayscale (quotes)}}\\[-1pt]
  {\scriptsize\color{acgray}2026-05-21}
\end{minipage}\hfill\begin{minipage}[t]{0.194\textwidth}\centering
  
\includegraphics[width=\linewidth,height=\linewidth,keepaspectratio]{figures/assets/c13-hellsine.png}\\[2pt]
  {\scriptsize\textbf{pop cover glyph}}\\[-1pt]
  {\scriptsize\color{acgray}2026-05-24}
\end{minipage}\hfill\begin{minipage}[t]{0.194\textwidth}\centering
  
\includegraphics[width=\linewidth,height=\linewidth,keepaspectratio]{figures/assets/c06-pink.png}\\[2pt]
  {\scriptsize\textbf{keymaps pink}}\\[-1pt]
  {\scriptsize\color{acgray}2026-06-13}
\end{minipage}\hfill\begin{minipage}[t]{0.194\textwidth}\centering
  
\includegraphics[width=\linewidth,height=\linewidth,keepaspectratio]{figures/assets/c07-blue.png}\\[2pt]
  {\scriptsize\textbf{keymaps blue}}\\[-1pt]
  {\scriptsize\color{acgray}2026-06-13}
\end{minipage}

\vspace{0.8em}

\begin{minipage}[t]{0.194\textwidth}\centering
  
\includegraphics[width=\linewidth,height=\linewidth,keepaspectratio]{figures/assets/c08-green.png}\\[2pt]
  {\scriptsize\textbf{keymaps green}}\\[-1pt]
  {\scriptsize\color{acgray}2026-06-13}
\end{minipage}\hfill\begin{minipage}[t]{0.194\textwidth}\centering
  
\includegraphics[width=\linewidth,height=\linewidth,keepaspectratio]{figures/assets/c09-orange.png}\\[2pt]
  {\scriptsize\textbf{keymaps orange}}\\[-1pt]
  {\scriptsize\color{acgray}2026-06-13}
\end{minipage}\hfill\begin{minipage}[t]{0.194\textwidth}\centering
  
\includegraphics[width=\linewidth,height=\linewidth,keepaspectratio]{figures/assets/c10-purple.png}\\[2pt]
  {\scriptsize\textbf{keymaps purple}}\\[-1pt]
  {\scriptsize\color{acgray}2026-06-13}
\end{minipage}\hfill\begin{minipage}[t]{0.194\textwidth}\centering
  
\includegraphics[width=\linewidth,height=\linewidth,keepaspectratio]{figures/assets/c11-yellow.png}\\[2pt]
  {\scriptsize\textbf{keymaps yellow}}\\[-1pt]
  {\scriptsize\color{acgray}2026-06-13}
\end{minipage}\hfill

\newpage
\section*{Appendiks C --- Live placeringer (rigtige UI'er)}
\addcontentsline{toc}{section}{Appendiks C --- Live placeringer}
\noindent \pals{} fotograferet levende i rigtige produktionsgrænseflader,
hver tekstet med sin kilde-URL (links er live): Visual Studio
Marketplace-opslaget (\pals{} som udvidelsens produktikon), monorepo-README'en
på både GitHub og den kanoniske Tangled-host og det roterende
\href{https://pals.aesthetic.computer}{\texttt{pals.aesthetic.computer}}-endpoint
(samme URL, forskellige positurer ved genindlæsning --- plus den rå
\texttt{/api/logo.png} givet til \texttt{<img>}-tags og avatarer).
\vspace{1.0em}

% Hand-authored: the Pals living in real production UIs, each with a caption + live URL.

\newcommand{\uishot}[3]{%
  % #1 image  #2 caption  #3 url
  \begin{center}
  \fbox{\includegraphics[width=0.96\linewidth]{#1}}\\[0.4em]
  {\footnotesize #2}\\[0.1em]
  {\footnotesize\ttfamily\href{#3}{#3}}
  \end{center}
  \vspace{0.7em}
}
\newcommand{\uishottwo}[6]{%
  % (#1 img #2 cap #3 url) (#4 img #5 cap #6 url)
  \noindent
  \begin{minipage}[t]{0.485\linewidth}\centering
    \fbox{\includegraphics[width=0.97\linewidth]{#1}}\\[0.35em]
    {\footnotesize #2}\\[0.1em]{\scriptsize\ttfamily\href{#3}{#3}}
  \end{minipage}\hfill
  \begin{minipage}[t]{0.485\linewidth}\centering
    \fbox{\includegraphics[width=0.97\linewidth]{#4}}\\[0.35em]
    {\footnotesize #5}\\[0.1em]{\scriptsize\ttfamily\href{#6}{#6}}
  \end{minipage}
  \vspace{0.9em}
}

\uishot{figures/live/ui-vscode-marketplace.png}{%
  \textbf{Visual Studio Marketplace} --- the \pals{} as the product icon of
  the official \ac{} VS~Code extension (1{,}201 installs).}{%
  https://marketplace.visualstudio.com/items?itemName=aesthetic-computer.aesthetic-computer-code}

\uishottwo
  {figures/live/ui-github-readme.png}{\textbf{GitHub} --- the \pals{} atop the
   monorepo README.}{https://github.com/whistlegraph/aesthetic-computer}
  {figures/live/ui-tangled-readme.png}{\textbf{Tangled} (canonical git host)
   --- the same README, same mark.}{https://tangled.org/aesthetic.computer/core}

\uishot{figures/compare/pals-live-strip.png}{%
  \textbf{\texttt{pals.aesthetic.computer}} in a live browser --- the same URL
  reloaded shows different poses (the rotating ``logo grabber''); the raw
  \texttt{/api/logo.png} is what \texttt{<img>} tags and avatars (Gravatar,
  social) fetch across the web.}{%
  https://pals.aesthetic.computer}



\end{document}
